%Data institusi

%\input{Dozentdaten}

%\renewcommand{\fachbereich}{Bidang studi}

%\renewcommand{\dozent}{Prof. Dr. . }

%Data ujian

\renewcommand{\klausurgebiet}{ }

\renewcommand{\klausurtyp}{ }

\renewcommand{\klausurdatum}{. 20}

\klausurvorspann {\fachbereich} {\klausurdatum} {\dozent} {\klausurgebiet} {\klausurtyp}

%Data untuk tabel poin berikut

\renewcommand{\aeins}{  3  }

\renewcommand{\azwei}{  3   }

\renewcommand{\adrei}{  5  }

\renewcommand{\avier}{  2  }

\renewcommand{\afuenf}{  4  }

\renewcommand{\asechs}{  8  }

\renewcommand{\asieben}{  4  }

\renewcommand{\aacht}{  5  }

\renewcommand{\aneun}{  3  }

\renewcommand{\azehn}{  2  }

\renewcommand{\aelf}{  6  }

\renewcommand{\azwoelf}{  7  }

\renewcommand{\adreizehn}{  4   }

\renewcommand{\avierzehn}{  8  }

\renewcommand{\afuenfzehn}{  64  }

\renewcommand{\asechzehn}{    }

\renewcommand{\asiebzehn}{    }

\renewcommand{\aachtzehn}{    }

\renewcommand{\aneunzehn}{    }

\renewcommand{\azwanzig}{    }

\renewcommand{\aeinundzwanzig}{    }

\renewcommand{\azweiundzwanzig}{    }

\renewcommand{\adreiundzwanzig}{    }

\renewcommand{\avierundzwanzig}{    }

\renewcommand{\afuenfundzwanzig}{    }

\renewcommand{\asechsundzwanzig}{    }

\punktetabellevierzehn

\klausurnote

\newpage

\setcounter{section}{0}

\inputaufgabegibtloesung
{3}
{

Definisikan istilah-istilah berikut
\zusatzklammer {yang dicetak miring} {} {}.
\aufzaehlungsechs{\stichwort {turunan tangensial kedua} {}
\zusatzklammer {atau
\stichwort {percepatan tangensial} {}} {} {}
\zusatzklammer {sebagai pemetaan ke $\R^n$} {} {}
dari suatu
\definitionsverweis {kurva yang dua kali terdiferensialkan secara kontinu}{}{}
\maabbdisp {\gamma} {I} {Y
} {}
yang bernilai dalam suatu
\definitionsverweis {hipermuka terdiferensialkan}{}{}

\mavergleichskette
{\vergleichskette
{ Y
}
{ \subseteq }{ \R^n
}
{  }{
}
{    }{
}
{   }{
}
}
{}{}{.}

}{Suatu ruang topologis $X$ yang
\stichwort {terhubung} {}.

}{\stichwort {Pemetaan singgung di titik} {}

\mavergleichskette
{\vergleichskette
{  P
}
{ \in }{ M
}
{  }{}
{    }{}
{   }{}
}
{}{}{}
untuk suatu pemetaan diferensiabel
\maabbdisp {\varphi} { M } { N
} {,}
dengan
\mathkor {} {M} {dan} {N} {}
manifold diferensiabel.

}{\stichwort {Produk} {}
dari manifold diferensiabel
\mathkor {} {M} {dan} {N} {.}

}{\stichwort {Ukuran volume} {}
yang terkait dengan bentuk volume positif $\omega$ pada suatu manifold diferensiabel berdimensi $n$.

}{Suatu
\stichwort {kurva geodesik} {}
pada manifold Riemann $M$.
}

}
{} {}

\inputaufgabegibtloesung
{3}
{

Nyatakan teorema-teorema berikut.
\aufzaehlungdrei{\stichwort {Teorema tentang kelengkungan kurva bidang yang diberikan secara implisit} {}

\mavergleichskette
{\vergleichskette
{  Y
}
{ \subseteq }{  \R^2
}
{  }{
}
{    }{
}
{   }{
}
}
{}{}{.}}{\stichwort {Rumus transformasi untuk bentuk volume positif pada manifold} {.}}{Teorema tentang kelengkungan Gauss dan kelengkungan seksional pada suatu permukaan diferensiabel

\mavergleichskette
{\vergleichskette
{  Y
}
{ \subseteq }{  \R^3
}
{  }{
}
{    }{
}
{   }{
}
}
{}{}{}}

}
{} {}

\inputaufgabegibtloesung
{5}
{

Misalkan

\mavergleichskette
{\vergleichskette
{ G
}
{ \subseteq }{ \R^n
}
{  }{
}
{    }{
}
{   }{
}
}
{}{}{}
\definitionsverweis {terbuka}{}{}
dan
\maabbdisp {\varphi} { G } { \R^m
} {}
adalah
\definitionsverweis {pemetaan yang terdiferensialkan secara kontinu}{}{,}
yang pada titik

\mavergleichskette
{\vergleichskette
{ P
}
{ \in }{ G
}
{  }{
}
{    }{
}
{   }{
}
}
{}{}{}
memiliki
\definitionsverweis {diferensial total}{}{}
yang surjektif. Misalkan

\mavergleichskette
{\vergleichskette
{ v
}
{ \in }{ T_PZ
}
{  }{
}
{    }{
}
{   }{
}
}
{}{}{}
adalah suatu vektor dalam
\definitionsverweis {ruang singgung serat}{}{}
$Z$ dari $\varphi$ yang melalui $P$. Tunjukkan bahwa terdapat
\definitionsverweis {kurva yang terdiferensialkan secara kontinu}{}{}
\maabbdisp {\gamma} { ]- \epsilon, \epsilon[ } { Z
} {}
\zusatzklammer {untuk suatu

\mavergleichskettek
{\vergleichskettek
{ \epsilon
}
{ > }{ 0
}
{  }{
}
{    }{
}
{   }{
}
}
{}{}{}
yang sesuai} {} {}
dengan

\mavergleichskette
{\vergleichskette
{ \gamma(0)
}
{ = }{ P
}
{  }{
}
{    }{
}
{   }{
}
}
{}{}{}
dan

\mavergleichskettedisp
{\vergleichskette
{ \gamma'(0)
}
{ =} { v
}
{ } {
}
{ } {
}
{ } {
}
}
{}{}{}.

}
{} {}

\inputaufgabegibtloesung
{2}
{

Misalkan
\mathkor {} {r} {dan} {R} {}
adalah bilangan real positif dengan

\mavergleichskette
{\vergleichskette
{ 0
}
{ < }{ r
}
{ < }{ R
}
{    }{
}
{   }{
}
}
{}{}{.}
Tentukan suatu
\definitionsverweis {medan normal satuan}{}{}
untuk torus
\zusatzklammer {tertanam} {} {}

\mavergleichskettedisp
{\vergleichskette
{ T
}
{ =} { \left\{ (x,y,z) \in \R^3  \mid   \left(  \sqrt{x^2+y^2} -R  \right)^2 +z^2  = r^2         \right\}
}
{ } {
}
{ } {
}
{ } {
}
}
{}{}{.}

}
{} {}

\inputaufgabe
{4}
{

Jelaskan peran teorema tentang pemetaan implisit dalam pengembangan teori manifold.

}
{} {}

\inputaufgabegibtloesung
{8}
{

Buktikan teorema tentang sifat adjoin-diri (self-adjoint) dari pemetaan Weingarten.

}
{} {}

\inputaufgabegibtloesung
{4}
{

Misalkan $X$ adalah suatu
\definitionsverweis {ruang kompak}{}{}
dan misalkan

\mavergleichskette
{\vergleichskette
{Y
}
{ \subseteq }{ X
}
{  }{
}
{    }{
}
{   }{
}
}
{}{}{}
adalah
\definitionsverweis {himpunan bagian tertutup}{}{,}
yang dilengkapi dengan
\definitionsverweis {topologi terinduksi}{}{.}
Tunjukkan bahwa $Y$ juga kompak.

}
{} {}

\inputaufgabegibtloesung
{5}
{

Berikan contoh suatu
\definitionsverweis {manifold diferensiabel}{}{}
$M$ yang
\definitionsverweis {berdimensi dua}{}{}
dan
\definitionsverweis {terhubung}{}{,}
beserta suatu titik

\mavergleichskette
{\vergleichskette
{ P
}
{ \in }{ M
}
{  }{
}
{    }{
}
{   }{
}
}
{}{}{}
sedemikian sehingga
\mathkor {} {M} {dan} {M \setminus \{P\}} {}
saling
\definitionsverweis {difeomorfik}{}{.}

}
{} {}

\inputaufgabegibtloesung
{3}
{

Untuk setiap
\definitionsverweis {vektor singgung}{}{}

\mavergleichskette
{\vergleichskette
{ v
}
{ \in }{ T_PS^2
}
{ \subseteq }{ \R^3
}
{    }{
}
{   }{
}
}
{}{}{}
dengan

\mavergleichskette
{\vergleichskette
{ \Vert {v} \Vert
}
{ = }{ 1
}
{  }{
}
{    }{
}
{   }{
}
}
{}{}{}
di suatu titik

\mavergleichskette
{\vergleichskette
{ P
}
{ \in }{ S^2
}
{ \subset }{ \R^3
}
{    }{
}
{   }{
}
}
{}{}{}
pada
\definitionsverweis {permukaan bola satuan}{}{,}
berikan suatu wakil
\definitionsverweis {diferensiabel}{}{}
\maabbdisp {\gamma} { I } { S^2
} {}
dengan

\mavergleichskette
{\vergleichskette
{ \gamma(0)
}
{ = }{ P
}
{  }{
}
{    }{
}
{   }{
}
}
{}{}{.}

}
{} {}

\inputaufgabegibtloesung
{2}
{

Tinjau parabola standar

\mavergleichskette
{\vergleichskette
{ Y
}
{ \subseteq }{ \R^2
}
{  }{
}
{    }{
}
{   }{
}
}
{}{}{}
dengan struktur Riemann terinduksi dan parametrisasi
\maabbeledisp {\varphi} {\R} { Y
} { t } { \left(  t   , \,  t^2                   \right)
} {,}
yang kita pandang sebagai suatu peta
\zusatzklammer {invers} {} {}.
Tentukan
\definitionsverweis {fungsi fundamental}{}{}
Riemann

\mavergleichskette
{\vergleichskette
{ g(t)
}
{ = }{ g_{11} (t)
}
{  }{
}
{    }{
}
{   }{
}
}
{}{}{.}

}
{} {}

\inputaufgabegibtloesung
{6}
{

Hitung bentuk diferensial tarik balik
\mathl{\varphi^* \tau}{} dari

\mavergleichskettedisp
{\vergleichskette
{ \tau
}
{ =} { dx \wedge dy \wedge dz - w dx \wedge dy \wedge dw + \cos (xy)  dx \wedge dz \wedge dw-yw dy \wedge dz \wedge dw
}
{ } {
}
{ } {
}
{ } {
}
}
{}{}{}
oleh pemetaan
\maabbeledisp {\varphi} { \R^3 } { \R^4
} { (r,s,t) } {  \left(  r^2s,t, \sin  r ,e^{st}  \right)  = (x,y,z,w)
} {.}

}
{} {}

\inputaufgabegibtloesung
{7 (2+3+2)}
{

Tinjau pemetaan-pemetaan diferensiabel
\maabbeledisp {\gamma} { [1,2] } {\R^2
} { t } {(t,t^{-1})
} {,}
dan
\maabbeledisp {\varphi} {\R^2} {\R^3
} {(u,v)} {(u^2,uv,-u+v^2)
} {,}
serta bentuk diferensial

\mavergleichskettedisp
{\vergleichskette
{ \omega
}
{ =} { xdx-zdy+dz
}
{ } {
}
{ } {
}
{ } {
}
}
{}{}{}
pada $\R^3$.

a) Hitung bentuk diferensial tarik balik
\mathl{\varphi^* \omega}{} pada $\R^2$.

b) Hitung integral lintasan bentuk diferensial
\mathl{\varphi^* \omega}{} sepanjang lintasan $\gamma$.

c) Hitung
\zusatzklammer {tanpa menggunakan hasil b)} {} {}
integral lintasan bentuk diferensial
\mathl{\omega}{} sepanjang lintasan $\varphi \circ \gamma$.

}
{} {}

\inputaufgabegibtloesung
{4}
{

Misalkan $M$ adalah suatu
\definitionsverweis {manifold berbatas}{}{.}
Tunjukkan bahwa setiap himpunan bagian terbuka

\mavergleichskette
{\vergleichskette
{ N
}
{ \subseteq }{ M
}
{  }{
}
{    }{
}
{   }{
}
}
{}{}{}
juga merupakan manifold dengan batas
\zusatzklammer {yang mungkin kosong} {} {},
dan bahwa

\mavergleichskettedisp
{\vergleichskette
{ \partial N
}
{ =} { \partial M \cap N
}
{ } {
}
{ } {
}
{ } {
}
}
{}{}{}
berlaku.

}
{} {}

\inputaufgabegibtloesung
{8}
{

Buktikan teorema tentang karakterisasi kurva geodesik terhadap koneksi Levi-Civita.

}
{} {}
